%définir l'entête de ce chapitre 
\markboth{Introduction Générale}{}

Le secteur de la gestion des ressources humaines connaît aujourd'hui une transformation profonde, portée par les avancées rapides des technologies numériques. Les organisations, qu'elles soient de grande envergure ou de taille intermédiaire, sont confrontées à des défis croissants en matière d'administration du personnel, de suivi des performances et de pilotage stratégique des ressources humaines. Dans ce contexte en perpétuelle mutation, la capacité à centraliser les données, à automatiser les processus et à intégrer des outils intelligents est devenue un facteur déterminant de compétitivité et d'efficacité organisationnelle.

Les PME et ETI souffrent particulièrement de l'absence de solutions RH à la fois complètes, flexibles et accessibles. Les systèmes existants sur le marché présentent soit une complexité technique et des coûts prohibitifs inadaptés aux petites structures, soit des fonctionnalités trop limitées pour répondre pleinement aux besoins métier. Face à cette lacune identifiée, une décision stratégique a été prise de concevoir et de développer une plateforme intégrée de gestion des ressources humaines dénommée SMART RH 4.0, capable d'offrir une expérience utilisateur optimale et de permettre aux organisations d'accéder facilement à une suite complète d'outils et de services RH.

C'est dans ce contexte que s'inscrit ce projet de fin d'études, réalisé au sein de l'entreprise dans le cadre de l'obtention du diplôme national d'ingénieur en génie logiciel. Ce stage représente également une opportunité de mettre en pratique les connaissances théoriques acquises tout au long du cursus académique, de consolider les compétences techniques et de se confronter aux réalités du monde professionnel en tant que développeur et architecte de solutions d'entreprise.

Notre mission consiste à concevoir et développer la plateforme SMART RH 4.0, en visant une expérience utilisateur fluide et intuitive, tout en permettant aux organisations de bénéficier d'une gamme étendue de services RH intégrés, couvrant la gestion des employés, le suivi des performances, la gestion de la paie ainsi que le recrutement.

Ce rapport retrace les différentes étapes de réalisation du projet. Il est structuré en quatre chapitres, répartis comme suit :

\begin{itemize}
    \item Le premier chapitre, intitulé « Cadre général du projet », présente l'organisme d'accueil, expose le contexte et les enjeux du projet, détaille la mission confiée, ainsi que la méthodologie de travail adoptée pour sa réalisation.
    
    \item Le deuxième chapitre, « Analyse et spécification des besoins », se penche sur les étapes cruciales de la démarche à suivre pour comprendre les exigences fonctionnelles et non fonctionnelles de la plateforme SMART RH 4.0. Il identifie également les acteurs clés impliqués et présente de manière claire les diagrammes de cas d'utilisation ainsi que le diagramme de classes, permettant une vision d'ensemble du système.
    
    \item Le troisième chapitre, « Aperçu conceptuel », approfondit la planification en détaillant les diagrammes de cas d'utilisation pour chaque fonctionnalité spécifique, les diagrammes de séquence, ainsi qu'une analyse de la conception architecturale et de la structure des données retenue pour notre solution.
    
    \item Enfin, le quatrième chapitre, « Réalisation », est consacré à l'étude technique du projet. Il spécifie les technologies et les outils utilisés, et présente la réalisation concrète des modules clés de la plateforme SMART RH 4.0, accompagnée de captures d'écran représentatives des interfaces développées tout au long de la période de stage.
\end{itemize}

Ce rapport se conclut par une synthèse récapitulant le travail accompli, les résultats obtenus, ainsi que les perspectives d'évolution et d'amélioration envisagées pour la plateforme SMART RH 4.0.
